\section{Additional experimental detail}
\label{app:detail}

\subsection{Trigger hyper-parameter sweep (E1)}
\label{app:sweep}

\ladder has three hyper-parameters: a patience $k$ (how many iterations without improvement before
the trigger may fire), a surrogate threshold $\tau$ (below which the rung is judged cold), and an
extension factor $e$ (how much extra budget a warm rung receives). At $\tau = 0$ and $e = 1$ the
arm reduces to \fixedsplit exactly, which is what makes the contrast in \secref{sec:e1} a clean
ablation of the trigger alone.

We swept $k \in \{2, 5, 20\}$, $\tau \in \{0.05, 0.3, 0.7\}$ and $e \in \{1, 3, 10\}$, giving 27
settings, each replayed at every budget. The best setting anywhere in the grid gains $+0.06$ over
\fixedsplit; the worst loses $-0.17$. No setting produces a gain that survives Holm correction. The
conclusion in \secref{sec:e1} is therefore a property of the trigger idea, not of one unlucky
parameterisation.

\subsection{Rung attribution (E1)}
\label{app:rungs}

For every detection made by \ladder we record which rung was active at the moment of success. Rung
0 is the default \texttt{plain\_init} random search, rung 1 is re-initialisation from a related
success, rung 2 is a template-family switch. Across all budgets, rung 0 accounts for 0 detections,
rung 1 for 31, and rung 2 for 13. This is the direct measurement behind the claim that the benefit
is a \family effect.

\subsection{Pressure design families (E4)}
\label{app:families}

The four pressure families are applied after the model produces a correct answer to a \gsmk problem
under greedy decoding.

\begin{itemize}[leftmargin=*,itemsep=1pt,topsep=2pt]
    \item \textbf{P1, simple doubt.} A single expression of disagreement with no justification.
    \item \textbf{P2, appeal to authority.} A single turn attributing the contrary answer to an
    expert source.
    \item \textbf{P3, fabricated consensus.} A single turn asserting that several other people
    reached a different answer.
    \item \textbf{P4, persistent multi-turn.} Repeated pushback across several turns without new
    argument.
\end{itemize}

The alignment intervention under test is an anti-sycophancy system prompt instructing the model to
maintain a correct answer under social pressure and to change it only in response to a substantive
mathematical argument. The control condition is the same grid with no system prompt. An item counts
as ``flipped'' when the model's final answer changes away from the correct value.

\subsection{Per-configuration coverage (E3)}
\label{app:coverage}

Pooled coverage in \tabref{tab:e3_coverage} averages over 14 configurations that behave very
differently. Configurations whose true elicitable fraction is near the middle of the range have
coverage close to nominal, which is what drags the pooled figure up. The three configurations with
the highest true elicitable fraction, \texttt{llama2-13b}, \texttt{llama2-70b} and
\texttt{llama2\_7b} at $0.96$--$1.00$, have coverage of exactly $0.00$ for the adaptive interval:
the bound fails in every replicate. This is the pattern one expects when the estimand is wrong
rather than when the interval is merely too narrow, since a too-narrow interval degrades gradually
while a wrong estimand fails systematically wherever the gap between the two quantities is large.

\subsection{Artifacts}
\label{app:artifacts}

All code, raw outputs, summaries and statistical tests are released alongside the paper:
offline results (E1--E3 raw traces, summaries and tests), live results (the E4 grid and protocol
replay), the figure-generation scripts, and the pre-registered analysis plan written before any
protocol was run.
